% !TeX root = ../../Book4.tex
% 纲要标题：### 6.4 Euler–Poincaré 等式与 Grothendieck 群
% 纲要标签：b4:sec:0504
% 纲要位置：计划纲要.md，第 3236 行。

\section{Euler–Poincaré 等式与 Grothendieck 群}
\label{b4:sec:0504}

边界对象在相邻次数中以相反符号出现，使复形各项的交错和等于同调的交错和。Grothendieck 群把这一计算推广到没有数值维数的交换范畴。

% 纲要内容（仅作写作计划，尚非教材正文）：
%
% * 有限长度或有限维情形的交错和
% * Grothendieck 群 K_0 中的可加不变量
% * 有界复形与其同调的 Euler 特征相等
%
% ---
%

\subsection{有限长度与有限维情形的交错和}
\label{b4:subsec:0504:01}

求出所有同调群有时很费力，但一个交错和却能在求核取商之前计算。原因是每个边界对象在两个相邻次数中各出现一次，且符号相反。

\begin{definition}\label{b4:def:euler-characteristic}
设 \(k\) 是域，\(C\) 是由有限维 \(k\)-向量空间组成的有界复形。规定
\[
\chi(C)\coloneqq\sum_n(-1)^n\dim_k C^n,
\]
称其为\term[Euler-tezheng]{Euler 特征}[Euler characteristic]。设 \(R\) 为含幺环；若 \(C\) 是由有限长度幺性左 \(R\) 模组成的有界上链复形，以 \(\ell_R(C^n)\) 表示第 \(n\) 项的合成长度，则把维数替换为此长度所得的有限交错和称为长度的 Euler 特征。
\end{definition}
\symidx{\(\chi(C)\)：有界有限维复形的 Euler 特征}
\symidx{\(\ell_R(M)\)：模的合成长度}

有界性保证只有有限多个非零项；每项有限维或有限长度，则保证每个被相加的数值都是有限整数。这是两个不同的要求。这里有限长度指具有有限合成列，其长度是合成因子的个数；长度的良定义与短正合列可加性已在第二卷命题12.2.5“长度的可加性”证明。有限维情形只需秩—零化度公式。两种情形都有
\[
0\to U\to V\to W\to0
\quad\Longrightarrow\quad
\lambda(V)=\lambda(U)+\lambda(W),
\]
其中 \(\lambda\) 分别表示维数或长度。因而对有界复形的短正合列，逐次相加即可得到 \(\chi(B)=\chi(A)+\chi(C)\)。

\begin{example}\label{b4:exm:graph-euler}
取一个有限图，顶点数为 \(v\)，边数为 \(e\)，连通分支数为 \(c\)。给每条边定向，在域 \(k\) 上定义链复形 \(C_1=k^e\xrightarrow{\partial}C_0=k^v\)，其余项为零。若用上链指标 \(C^{-n}=C_n\)，则 \((-1)^{-n}=(-1)^n\)，所以其 Euler 特征仍为 \(v-e\)。以 \(\varepsilon_w\) 表示顶点 \(w\) 的基向量，把一条边送到其终点与起点的基向量之差。每条边界在所属连通分支中的坐标和为零，故像包含在各分支的零和子空间中。反过来，在一个分支中选定基准顶点 \(w_0\)，从 \(w_0\) 到任意顶点 \(w\) 的路径按行进方向给各边加上正负号，其边界就是 \(\varepsilon_w-\varepsilon_{w_0}\)。这些差生成该分支的全部零和向量：若 \(\sum_w a_w=0\)，则
\[
\sum_w a_w\varepsilon_w
=\sum_{w\ne w_0}a_w(\varepsilon_w-\varepsilon_{w_0}).
\]
因此两个子空间相等，每个分支贡献的秩是其顶点数减一，总计 \(\operatorname{rank}\partial=v-c\)，从而
\[
\dim_kH_0=c,\qquad \dim_kH_1=e-v+c,
\qquad v-e=\dim_kH_0-\dim_kH_1.
\]
例如三角形的 \(H_1\) 为一维，而加上一条没有构成新回路的悬挂边不改变这个维数。
\end{example}

\subsection{\texorpdfstring{Grothendieck 群 \(K_0\) 中的可加不变量}{Grothendieck 群 K0 中的可加不变量}}
\label{b4:subsec:0504:02}

维数把所有一维向量空间看成同一种贡献；一般交换范畴中可能有几种不能互相替代的简单对象。为了保留这种差别，可以只规定短正合列应满足的加法关系。Euler 交错和还需要取对象类的差，因此不能只在交换幺半群中相加；下面从自由交换群出发，允许有限整数线性组合，包括负系数。

\begin{definition}\label{b4:def:grothendieck-group}
设 \(\mathcal A\) 是本质小的交换范畴，即其对象同构类组成一个集合。取以各同构类 \([M]\) 为基的自由交换群，再除去全部关系
\[
[B]-[A]-[C]
\qquad\bigl(0\to A\to B\to C\to0\;\text{短正合}\bigr)
\]
所生成的子群，所得群记为 \(K_0(\mathcal A)\)，称为 \(\mathcal A\) 的\term[Grothendieck-qun]{Grothendieck 群}[Grothendieck group]。
\end{definition}
\symidx{\(K_0(\mathcal A)\)：交换范畴的 Grothendieck 群}

本质小的条件保证上述自由群以集合为基。对于全部模组成的大范畴，需要先限制大小或选取所关心的本质小子范畴；不能不加说明地对真类取自由群。本节的交换范畴版本对全部短正合列施加关系。环的 \(K_0(R)\) 则以有限生成投射模的同构类为生成元；它们之间态射在模范畴中的核、余核未必仍为有限生成投射模，因而此范畴通常不是交换范畴。三项都是有限生成投射模的短正合列必因商模投射而分裂，所以这里使用的关系是 \([A\oplus C]=[A]+[C]\)。对象范围和所用的正合列均须区分，不能仅凭同一个符号将两种群等同。\Cref{b4:def:B-grothendieck-group}会以指定正合结构统一这两种定义；\cref{b4:thm:B-euler-isomorphism}进一步证明投射模与完美复形的 Euler 类给出相同的 Grothendieck 群。

\begin{proposition}\label{b4:prop:grothendieck-universal}
设 \(\mathcal A\) 是本质小交换范畴，\(G\) 是交换群。对 \(\mathcal A\) 的对象同构类赋值 \(\lambda(M)\in G\)，若每条短正合列 \(0\to A\to B\to C\to0\) 均满足 \(\lambda(B)=\lambda(A)+\lambda(C)\)，则存在唯一群同态 \(\bar\lambda:K_0(\mathcal A)\to G\)，使 \(\bar\lambda([M])=\lambda(M)\)。若 \(\mathcal B\) 也是本质小交换范畴，则每个正合函子 \(F:\mathcal A\to\mathcal B\) 诱导群同态 \(K_0(F):K_0(\mathcal A)\to K_0(\mathcal B)\)，\([M]\mapsto[F(M)]\)。
\end{proposition}
\begin{proof}
由自由交换群的泛性质，生成元上的赋值唯一延拓为群同态。可加性恰好说定义 \(K_0\) 的每个关系都映到零，故唯一地通过商群分解。第二项中正合函子把定义关系送到定义关系，因此同样可以下降。关系 \([0]=[0]+[0]\) 还说明 \([0]=0\)，不必另外添加零对象关系。
\end{proof}

\begin{example}\label{b4:exm:kzero-vectors}
设 \(k\) 为域。有限维 \(k\)-向量空间范畴满足 \([V]=(\dim_kV)[k]\)，因为基给出 \(V\cong k^{\dim V}\)，且直和对应分裂短正合列。维数诱导 \(K_0\to\mathbb Z\)，而 \(1\mapsto[k]\) 是其逆，故 \(K_0\cong\mathbb Z\)。若改为有限维 \(k\times k\)-模，令 \(e_1=(1,0)\)、\(e_2=(0,1)\)，则 \(M=e_1M\oplus e_2M\)。记两种简单模为 \(S_1=k\times0\)、\(S_2=0\times k\)，并令 \(a=\dim_k e_1M\)、\(b=\dim_k e_2M\)。在各分量中选基给出
\[
e_1M\cong S_1^{\oplus a},\qquad
e_2M\cong S_2^{\oplus b},\qquad
[M]=a[S_1]+b[S_2].
\]
模同态保持这两个分量，短正合列逐分量正合，所以两坐标维数诱导 \(K_0\to\mathbb Z^2\)。它与群同态 \((a,b)\mapsto a[S_1]+b[S_2]\) 互为逆：两种简单模分别映到 \((1,0)\)、\((0,1)\)，而上式把任意 \([M]\) 恢复出来。因此 \(K_0\cong\mathbb Z^2\)。总维数只是这两个坐标之和，不能区别两种简单模。
\end{example}

\subsection{有界复形与同调的 Euler 特征}
\label{b4:subsec:0504:03}

\begin{theorem}[Euler–Poincaré 等式]\label{b4:thm:euler-poincare}
设 \(\mathcal A\) 是本质小交换范畴，\(C\) 是其中的有界复形。则在 \(K_0(\mathcal A)\) 中有
\[
\chi_{K_0}(C)\coloneqq\sum_n(-1)^n[C^n]
=\sum_n(-1)^n[H^n(C)].
\]
这称为\term[Euler-Poincare-dengshi]{Euler–Poincaré 等式}[Euler–Poincaré formula]。因此任何取值于交换群的短正合列可加不变量，都使复形与其同调的交错和相等。
\end{theorem}
\symidx{\(\chi_{K_0}(C)\)：取值于 Grothendieck 群的 Euler 类}
\begin{proof}
将微分分解为到像的满态射再接像的包含态射：
\[
\mathrm{d}^n:C^n\xrightarrow{e^n}B^{n+1}(C)
\xrightarrow{i^{n+1}}C^{n+1},
\qquad \mathrm{d}^n=i^{n+1}e^n.
\]
因为 \(i^{n+1}\) 为单态射，\(\ker e^n=\ker\mathrm{d}^n=Z^n(C)\)。因此对每个 \(n\)，有两条短正合列
\[
0\to Z^n(C)\to C^n\xrightarrow{e^n}B^{n+1}(C)\to0,
\qquad
0\to B^n(C)\to Z^n(C)\to H^n(C)\to0.
\]
第二条来自同调的定义。由这两条列得
\[
[C^n]=[B^n(C)]+[H^n(C)]+[B^{n+1}(C)].
\]
取整数 \(a\le b\)，使 \(C^n=0\) 对 \(n<a\) 或 \(n>b\) 成立。于是 \(B^a(C)=\operatorname{im}\mathrm{d}^{a-1}=0\)，且 \(B^{b+1}(C)=\operatorname{im}\mathrm{d}^b=0\)。将上式乘 \((-1)^n\) 并从 \(a\) 到 \(b\) 求和，其中一组边界项重新编号为
\[
\begin{aligned}
\sum_{n=a}^{b}(-1)^n[B^{n+1}(C)]
&=-\sum_{j=a+1}^{b+1}(-1)^j[B^j(C)]\\
&=-\sum_{j=a}^{b}(-1)^j[B^j(C)].
\end{aligned}
\]
最后一个等号用了上述两个端点为零。这组有限和恰与另一组边界项相消，剩下的正是所需公式。最后对该等式应用\cref{b4:prop:grothendieck-universal}的群同态即可。
\end{proof}

\begin{corollary}\label{b4:cor:euler-quasi-invariance}
设 \(k\) 为域、\(R\) 为含幺环。对有限维 \(k\) 向量空间的有界上链复形，或有限长度幺性左 \(R\) 模的有界上链复形，复形各项的维数或长度交错和等于其上同调的相应交错和。若同一类中的复形映射 \(f:C\to D\) 在每个次数诱导上同调同构，则 \(C,D\) 的相应 Euler 特征相等；此类有界零调复形的相应 Euler 特征为零。
\end{corollary}
\begin{proof}
边界、余圈及上同调仍为有限维向量空间或有限长度模，所以刚才的两条短正合列中，各项的维数或长度均有定义。以 \(\lambda\) 表示这一数值，由可加性得
\[
\lambda(C^n)=\lambda(B^n(C))+\lambda(H^n(C))+\lambda(B^{n+1}(C)).
\]
按前一证明的有限重新编号消去边界项，便得到所需等式。上同调同构使等式右端逐项相同；零调时每项为零。
\end{proof}

有限性不能只当作便于书写的条件。例如令 \(C^n=k\)（\(n\ge0\)）、\(C^n=0\)（\(n<0\)），并取 \(\mathrm{d}^{2r}=1_k\)、\(\mathrm{d}^{2r+1}=0\)（\(r\ge0\)），即
\[
0\to k\xrightarrow{1}k\xrightarrow{0}k\xrightarrow{1}k\xrightarrow{0}\cdots,
\]
其中第一个 \(k\) 位于零次。这是零调复形，但 \(1-1+1-1+\cdots\) 不是这里定义的有限交错和。\(K_0(\mathcal A)\) 的群运算只给出有限求和，并未定义无穷级数运算，所以不能将相邻项成对抵消后宣称这个复形的 Euler 类为零。若要给无界复形赋予更广的 Euler 特征，必须另给允许无限求和的结构和条件，或使所求交错和仍只有有限个非零项；本定理并不提供这样的延拓。
